\scanpage{35}
% archive sheet number: 34
\begin{center}\uline{Quasi-fonctions} \uline{de Néron-Weil}\end{center}

$K$ corps muni d'un ens ``propre'' $M_K$ de val. abs. qu'on écrit sous forme logarithmique $v(x)$

$\overline K$ une clôture \deleted{algébrique} \inserted{\deleted{séparable}}, $M_{\overline K}$ l'ens des val. absolues sur $\overline K$ prolongeant celles de $K$. Si $\Gamma$ est le gpe de galois de $\overline K$ sur $K$, on a
\[
M_{\overline K}/\Gamma\simeq M_K
\]
\margin{$\circledast$}

\uline{Fonctions réelles sur} $X(\overline K)\times M_{\overline K}$, \deleted{\unsure{quasi}} \uline{rationalité}.

Soit $X$ un préschéma sur $K$, \deleted{[\unsure{quelconque}} \inserted{et soit}
\[
\varphi:X(\overline K)\times M_{\overline K}\longrightarrow\mathbb R
\]
une fonction, notée $\varphi(P,v)$. Supposons que $\varphi$ commute à $\Gamma$ :
\[
\varphi(\sigma P,\sigma v)=\varphi(P,v)
\]
alors \deleted{pour tout} si \deleted{\unsure{$\sigma$}} $P\in X(\overline K)^\Gamma=X(K^{p^{-\infty}})$ (où $K^{p^{-\infty}}$ est la clôture parfaite de $K$) \deleted{\unsure{alors}} on aura $\varphi(\deleted{\sigma}P,\sigma v)=\varphi(P,v)$ pour tt $v,\sigma$, donc en vertu de $\circledast$ $\varphi(P,v)$ est connu quand on \unsure{connait|connaît} \deleted{\unsure{$\sigma$}} $v|K$.

Donc $\varphi$ induit
\[
\deleted{\unsure{A|X}}X(K^{p^{-\infty}})\times M_{\unsure{?|2}K}\longrightarrow\mathbb R
\]
\unsure{à|:} \unsure{fortiori|fonction} $X(K)\times M_K\longrightarrow\mathbb R$
